%% paginas/3_textuais/conclusao.tex
%% Elemento Textual Obrigatório (Conclusão) -- Conforme ABNT NBR 14724 & NBR 6024

\chapter{Conclusão}

A conclusão constitui a etapa final do documento acadêmico, onde o pesquisador realiza um fechamento integrativo sobre a sua investigação. Nela, retoma-se sinteticamente o problema de pesquisa, respondendo se os objetivos gerais e específicos foram integralmente atingidos, expondo as principais contribuições intelectuais do trabalho e sugerindo rumos e aperfeiçoamentos para futuros pesquisadores na mesma área temática.

No contexto do desenvolvimento deste ecossistema, conclui-se que a modularização de arquivos LaTeX voltados às normas ABNT é uma solução técnica viável e altamente recomendada. Ela remove a complexidade mecânica da diagramação dos olhos do pesquisador e promove uma barreira física saudável entre a lógica de estilo (\texttt{styles.sty}) e o conteúdo acadêmico distribuído. A criação de layouts limpos e dinâmicos para a Capa, Folha de Rosto e Aprovação permite que o modelo seja adaptado para qualquer instituição brasileira simplesmente alterando-se as constantes declaradas no arquivo central de dados (\texttt{metadata.tex}).

Como sugestões para trabalhos futuros, recomenda-se a expansão deste modelo modular com a integração direta a bibliotecas automatizadas como o \texttt{biblatex} (estilo ABNT), automatizações de scripts para conversão de planilhas locais para tabelas do padrão IBGE, e o desenvolvimento de macros para compilação contínua (CI/CD) gerando múltiplos formatos de saída digital e impressa automaticamente.
